\documentclass[10pt,a4paper]{amsart}
\usepackage[margin=19mm]{geometry}
\usepackage{amsmath,amssymb,mathtools}
\usepackage[scr=rsfs,cal=euler]{mathalfa}
\usepackage{xfrac}
\usepackage[unicode,hidelinks,bookmarks=false]{hyperref}
\newcommand{\ie}{i.e.\ }
\newcommand{\cf}{cf.\ }
\newcommand{\resp}{respectively\ }
\newcommand{\Subsection}{\subsection}
\providecommand{\nobreakdash}{\nobreak\hbox{-}}
\setlength{\emergencystretch}{2em}
\multlinegap0pt
\usepackage{bm}
\usepackage{bbm}
\usepackage{dsfont}
\usepackage{xspace}
\usepackage{cancel}
\usepackage{mathtools}
\usepackage{amsthm}
\makeatletter
\@ifundefined{thm}{\newtheorem{thm}{Thm}}{}
\@ifundefined{lem}{\newtheorem{lem}[thm]{Lemma}}{}
\@ifundefined{rem}{\newtheorem{rem}[thm]{Remark}}{}
\makeatother
\title[Homogeneous fractional integral operators on weighted Lebesg]{Homogeneous fractional integral operators on weighted Lebesgue, Morrey and Campanato spaces}
\author{Jingliang Du, Hua Wang}
\date{Source: arXiv:2510.00426v1 (CC BY 4.0)}
\begin{document}
\maketitle

\noindent\emph{Source and licence.} Homogeneous fractional integral operators on weighted Lebesgue, Morrey and Campanato spaces. Version arXiv:2510.00426v1 (CC BY 4.0), distributed under the Creative Commons Attribution 4.0 International licence, as recorded in the arXiv licence metadata for this exact version and shown on the abstract page https://arxiv.org/abs/2510.00426v1. Full rights notice: licenses/arxiv-2510.00426-CC-BY-4.0.txt.\par\smallskip

\noindent\textit{Editorial scope.} Counted unit: the Lem whose source label is \texttt{lem3} in the article (source0.tex lines 312--326, source label \texttt{lem3}), delivered together with the article's own complete original proof of it (source0.tex lines 327--393, one \texttt{proof} environment of 67 lines). No mathematics is written, repaired or re-derived: statement and proof are the article's own text line for line. Required context retained uncounted: 964 (lines 207--209); 965 (lines 219--221); remark10 (lines 225--231); lem1 (lines 298--309). Every definition, display and label of those lines is retained unchanged. Omitted from this package and not redistributed: 1--206, 210--218, 222--224, 232--297, 310--311, 394--858. Nothing inside a retained excerpt is omitted or altered: every retained body is delivered exactly as the article's lines are written, so no omission receipt of any kind accompanies this package. Citations: the retained excerpts carry no citation command at all, so no bibliography entry of the article is transcribed and no reference is redistributed. Reference adaptation: none was needed and none was made; every reference inside the delivered unit resolves inside it, so no unresolved reference or citation is produced. Printed slips kept literal: no printed word, symbol, hypothesis or number of the counted statement or of its proof was altered. Layout and class: the article's own class is \texttt{elsarticle} with packages including amsmath, amssymb, amsthm; the wrapper is the lane's pinned \texttt{amsart} editorial wrapper at 10pt on A4, which restates the theorem-like environments the retained text needs and adds the article's own macro definitions that the retained text needs (none), with extra wrapper packages (bm, bbm, dsfont, xspace, cancel, mathtools). The delivered numbering is the unit's own. Assets: no retained span contains a figure environment, a graphics-inclusion command, a PDF-page inclusion, a table or a TikZ picture; no image file is copied into this package, the package declares no asset dependency and no image file is redistributed with it. Licence: version arXiv:2510.00426v1 of this article is distributed under the Creative Commons Attribution 4.0 International licence, as recorded in the arXiv licence metadata for this exact version and shown on the abstract page https://arxiv.org/abs/2510.00426v1. A full rights notice is delivered at licenses/arxiv-2510.00426-CC-BY-4.0.txt. Characters: every retained excerpt is pure ASCII, so the delivered engine, XeLaTeX with its default Latin Modern text font, reports no missing character.
\begin{equation}\label{964}
\omega_s(\delta):=\sup_{|\rho|<\delta}\bigg(\int_{\mathbb{S}^{n-1}}\big|\Omega(\rho x')-\Omega(x')\big|^sd\sigma(x')\bigg)^{1/s},
\end{equation}

\begin{equation}\label{965}
\omega_{\infty}(\delta):=\sup_{|\rho|<\delta,x'\in \mathbb{S}^{n-1}}\big|\Omega(\rho x')-\Omega(x')\big|.
\end{equation}

\begin{rem}\label{remark10}
Let $1\leq s\leq\infty$. Since
\begin{equation*}
\int_0^1\frac{\omega_s(\delta)}{\delta}\,d\delta<\int_0^1\frac{\omega_s(\delta)}{\delta^{1+\beta}}\,d\delta
\end{equation*}
holds for any $\beta>0$, we know that if $\Omega$ satisfies the smoothness condition $\mathfrak{D}_{s;\beta}$, then it must satisfy the smoothness condition $\mathfrak{D}_{s}$.
\end{rem}

\begin{lem}\label{lem1}
Let $0<\alpha<n$. Suppose that $\Omega$ satisfies the $L^s$-Dini smoothness condition $\mathfrak{D}_{s}$ with $1<s\le\infty$. If there exists a constant $0<\vartheta<1/2$ such that if $|x|<\vartheta \mathcal{R}$, then we have
\begin{equation*}
\begin{split}
&\bigg(\int_{\mathcal{R}\leq|z|<2\mathcal{R}}
\bigg|\frac{\Omega(z-x)}{|z-x|^{n-\alpha}}-\frac{\Omega(z)}{|z|^{n-\alpha}}\bigg|^sdz\bigg)^{1/s}\\
&\leq C\cdot \mathcal{R}^{n/s-(n-\alpha)}\bigg(\frac{|x|}{\mathcal{R}}
+\int_{|x|/{2\mathcal{R}}}^{|x|/\mathcal{R}}\frac{\omega_{s}(\delta)}{\delta}d\delta\bigg),
\end{split}
\end{equation*}
where the constant $C>0$ is independent of $\mathcal{R}$ and $x$.
\end{lem}

\begin{lem}\label{lem3}
Let $0<\alpha<n$. Suppose that $\Omega$ satisfies the $L^{s;\beta}$-Dini smoothness condition $\mathfrak{D}_{s;\beta}$ with $1<s\le\infty$ and $0<\beta\leq1$. Then for any ball $\mathcal{B}=B(x_0,r)$ and for any $x,y\in \mathcal{B}$,
\begin{equation}\label{wewant0}
\begin{split}
&\bigg(\int_{2^{j+1}\mathcal{B}\backslash 2^j\mathcal{B}}
\left|\frac{\Omega(x-z)}{|x-z|^{n-\alpha}}-\frac{\Omega(y-z)}{|y-z|^{n-\alpha}}\right|^sdz\bigg)^{1/s}\\
&\leq C\cdot\big(2^jr\big)^{n/s-(n-\alpha)}
\bigg[\frac{1}{2^{j-1}}+\frac{1}{2^{j\beta}}
\int_{|x-x_0|/{(2^{j+1}r)}}^{|x-x_0|/{(2^jr)}}\frac{\omega_s(\delta)}{\delta^{1+\beta}}d\delta
+\frac{1}{2^{j\beta}}
\int_{|y-x_0|/{(2^{j+1}r)}}^{|y-x_0|/{(2^jr)}}\frac{\omega_s(\delta)}{\delta^{1+\beta}}d\delta\bigg].
\end{split}
\end{equation}
Here $2\leq j\in \mathbb{N}$ and $\omega_s$ is as in \eqref{964} or \eqref{965}.
\end{lem}

\begin{proof}
For any ball $\mathcal{B}=B(x_0,r)$ and $0<\alpha<n$, since
\begin{equation*}
\begin{split}
\left|\frac{\Omega(x-z)}{|x-z|^{n-\alpha}}-\frac{\Omega(y-z)}{|y-z|^{n-\alpha}}\right|
&\leq\left|\frac{\Omega(x-z)}{|x-z|^{n-\alpha}}-\frac{\Omega(x_0-z)}{|x_0-z|^{n-\alpha}}\right|\\
&+\left|\frac{\Omega(x_0-z)}{|x_0-z|^{n-\alpha}}-\frac{\Omega(y-z)}{|y-z|^{n-\alpha}}\right|,
\end{split}
\end{equation*}
it then follows from Minkowski's inequality that for each $j\geq2$,
\begin{equation*}
\begin{split}
&\bigg(\int_{2^{j+1}\mathcal{B}\backslash 2^j\mathcal{B}}
\left|\frac{\Omega(x-z)}{|x-z|^{n-\alpha}}-\frac{\Omega(y-z)}{|y-z|^{n-\alpha}}\right|^sdz\bigg)^{1/s}\\
\leq&
\bigg(\int_{2^{j+1}\mathcal{B}\backslash 2^j\mathcal{B}}\left|\frac{\Omega(x-z)}{|x-z|^{n-\alpha}}-\frac{\Omega(x_0-z)}{|x_0-z|^{n-\alpha}}\right|^sdz\bigg)^{1/s}\\
&+\bigg(\int_{2^{j+1}\mathcal{B}\backslash 2^j\mathcal{B}}
\left|\frac{\Omega(x_0-z)}{|x_0-z|^{n-\alpha}}-\frac{\Omega(y-z)}{|y-z|^{n-\alpha}}\right|^sdz\bigg)^{1/s}.
\end{split}
\end{equation*}
For any given $x\in\mathcal{B}$ and $z\in 2^{j+1}\mathcal{B}\backslash 2^j\mathcal{B}$ with $j\geq2$, it is easy to verify that
\begin{equation*}
|x-x_0|<\frac{\,1\,}{4}|z-x_0|.
\end{equation*}
As pointed out in Remark \ref{remark10}, for any $0<\beta\leq 1$ and $1<s\le\infty$, the condition $\mathfrak{D}_{s;\beta}$ is stronger than the $L^s$-Dini condition $\mathfrak{D}_{s}$. According to Lemma \ref{lem1}(here we take $\mathcal{R}=2^jr$), we get
\begin{equation*}
\begin{split}
&\bigg(\int_{2^{j+1}\mathcal{B}\backslash2^j\mathcal{B}}
\left|\frac{\Omega(x-z)}{|x-z|^{n-\alpha}}-\frac{\Omega(x_0-z)}{|x_0-z|^{n-\alpha}}\right|^sdz\bigg)^{1/s}\\
&=\bigg(\int_{2^jr\leq|z-x_0|<2^{j+1}r}\left|\frac{\Omega(z-x_0-(x-x_0))}{|z-x_0-(x-x_0)|^{n-\alpha}}
-\frac{\Omega(z-x_0)}{|z-x_0|^{n-\alpha}}\right|^sdz\bigg)^{1/s}\\
&\leq C\cdot\big(2^jr\big)^{n/s-(n-\alpha)}\bigg[\frac{|x-x_0|}{2^jr}+\int_{|x-x_0|/{(2^{j+1}r)}}^{|x-x_0|/{(2^jr)}}
\frac{\omega_s(\delta)}{\delta}d\delta\bigg].
\end{split}
\end{equation*}
Observe that for any $0<\beta\leq1$,
\begin{equation*}
\int_{|x-x_0|/{(2^{j+1}r)}}^{|x-x_0|/{(2^jr)}}\frac{\omega_s(\delta)}{\delta}d\delta
\leq\Big(\frac{|x-x_0|}{2^{j}r}\Big)^{\beta}
\int_{|x-x_0|/{(2^{j+1}r)}}^{|x-x_0|/{(2^jr)}}\frac{\omega_s(\delta)}{\delta^{1+\beta}}d\delta.
\end{equation*}
Hence, for any given $x\in \mathcal{B}=B(x_0,r)$,
\begin{equation}\label{wewant1}
\begin{split}
&\bigg(\int_{2^{j+1}\mathcal{B}\backslash2^j\mathcal{B}}
\left|\frac{\Omega(x-z)}{|x-z|^{n-\alpha}}-\frac{\Omega(x_0-z)}{|x_0-z|^{n-\alpha}}\right|^sdz\bigg)^{1/s}\\
&\leq C\cdot\big(2^jr\big)^{n/s-(n-\alpha)}\bigg[\frac{|x-x_0|}{2^jr}+\Big(\frac{|x-x_0|}{2^{j}r}\Big)^{\beta}
\int_{|x-x_0|/{(2^{j+1}r)}}^{|x-x_0|/{(2^jr)}}\frac{\omega_s(\delta)}{\delta^{1+\beta}}d\delta\bigg]\\
&\leq C\cdot\big(2^jr\big)^{n/s-(n-\alpha)}\bigg[\frac{1}{2^j}+\frac{1}{2^{j\beta}}
\int_{|x-x_0|/{(2^{j+1}r)}}^{|x-x_0|/{(2^jr)}}\frac{\omega_s(\delta)}{\delta^{1+\beta}}d\delta\bigg].
\end{split}
\end{equation}
Similarly, for any given $y\in \mathcal{B}=B(x_0,r)$, we can also obtain
\begin{equation}\label{wewant2}
\begin{split}
&\bigg(\int_{2^{j+1}\mathcal{B}\backslash 2^j\mathcal{B}}
\left|\frac{\Omega(x_0-z)}{|x_0-z|^{n-\alpha}}-\frac{\Omega(y-z)}{|y-z|^{n-\alpha}}\right|^sdz\bigg)^{1/s}\\
&\leq C\cdot\big(2^jr\big)^{n/s-(n-\alpha)}
\bigg[\frac{|y-x_0|}{2^jr}+\Big(\frac{|y-x_0|}{2^{j}r}\Big)^{\beta}
\int_{|y-x_0|/{(2^{j+1}r)}}^{|y-x_0|/{(2^jr)}}\frac{\omega_s(\delta)}{\delta^{1+\beta}}d\delta\bigg]\\
&\leq C\cdot\big(2^jr\big)^{n/s-(n-\alpha)}
\bigg[\frac{1}{2^j}+\frac{1}{2^{j\beta}}
\int_{|y-x_0|/{(2^{j+1}r)}}^{|y-x_0|/{(2^jr)}}\frac{\omega_s(\delta)}{\delta^{1+\beta}}d\delta\bigg].
\end{split}
\end{equation}
Combining the above estimates \eqref{wewant1} and \eqref{wewant2} leads to the estimate \eqref{wewant0}, which completes the proof of Lemma \ref{lem3}.
\end{proof}

\end{document}
